\documentclass[preprint,12pt]{elsarticle}

% Same fonts and macros as the manuscript, so a symbol means the same thing
% in both files. The supplement is uploaded as a PDF built here, not compiled
% by the journal's system, so it can refer back into the manuscript with xr:
% every \ref{M-...} below is resolved from theory.aux and fails the build if
% the manuscript's numbering moves. The manuscript refers forward to this
% file only in plain text ("Supplementary Section S1"), because the journal
% DOES compile the manuscript and would not have this file's .aux; those plain
% references are checked by verify_numbers.py instead.
%
% Sections S2 to S4 were moved here verbatim from theory.tex, once, by a
% script that cut each block by anchors matching exactly once and
% re-pointed its references at the manuscript. The move cannot be repeated
% -- the blocks are no longer in theory.tex -- so this file is edited
% directly from now on.
\usepackage{lmodern}
\usepackage[T1]{fontenc}
\usepackage{amsmath,amssymb,amsfonts}
\usepackage{amsthm}
\usepackage{array}
\usepackage{url}
\usepackage{xr}
\externaldocument[M-]{theory}

\tolerance=1500
\emergencystretch=2em
\hyphenpenalty=100
\hbadness=4000

\newcommand{\R}{\mathbb{R}}
\newcommand{\E}{\mathbb{E}}
\newcommand{\dvec}{\mathbf{d}}
\newcommand{\dhat}{\hat{\mathbf{d}}}
\newcommand{\Var}{\mathrm{Var}}
\newcommand{\Std}{\mathrm{Std}}
\newcommand{\dFR}{d_{\mathrm{FR}}}
\newcommand{\DB}{D_{\mathrm{B}}}
\newcommand{\lloss}{\ell}

\renewcommand{\thesection}{S\arabic{section}}
\renewcommand{\theequation}{S\arabic{equation}}
\renewcommand{\thetable}{S\arabic{table}}

\journal{Reliability Engineering \& System Safety}

\begin{document}

\begin{frontmatter}
\title{Supplementary material for: Information loss of fixed linear health
indicators as a control problem: simplex geometry and reachability limits}

% copied verbatim from theory.tex; make_supplement.py checks they still agree
\author[a]{Huy Hoang Le}
\ead{lhhoang@powermore.vn}
\author[b]{Kim-Anh Nguyen\corref{cor}}
\ead{nkanh@dut.udn.vn}
\cortext[cor]{Corresponding author}
\affiliation[a]{organization={R\&D Department, PowerMore Ltd.},
  city={Da Nang}, postcode={550000}, country={Viet Nam}}
\affiliation[b]{organization={Faculty of Electrical Engineering, University
  of Science and Technology, The University of Danang},
  city={Da Nang}, postcode={550000}, country={Viet Nam}}

\begin{abstract}
Section~\ref{sec:protocols} gives the protocol behind every row of
Table~\ref{M-tab:fleets} of the main text, with the parameters of the code
that produced it. Sections~\ref{sec:proofs} to~\ref{sec:verification} hold
material moved out of the main text: two proofs, the two candidate
certificates for Proposition~\ref{M-prop:converse} and why neither qualifies,
and the numerical verification of the identities and of the exceptional
family of Remark~\ref{M-rem:family}. Equation, section and proposition numbers
without an S refer to the main text.
\end{abstract}
\end{frontmatter}

\section{Protocols for the public-fleet tests}
\label{sec:protocols}

Every number in Table~\ref{M-tab:fleets} is recomputed by
\texttt{verify\_numbers.py} in the repository named in the main text, and the
parameters below are the ones that script passes; where a test was also run
at other settings to check stability, both are given. All four protocols read
an empirical informative direction as a time derivative of the monitoring
channels along normalised life, and so rely on the time-shift population of
Section~\ref{M-sec:prelim}, under which that derivative is the sensitivity of
Eq.~\eqref{M-eq:sens}.

\subsection{The rotation test on the Severson cells}

Channels are capacity at discharge, internal resistance and mean cell
temperature. A cell is used if its life is recorded, it has at least $400$
cycles, and all three channels are finite throughout; $129$ cells qualify.
Life is normalised per cell to $\tau\in[0,1]$ by cycle index. Each channel is
differentiated with a Savitzky--Golay filter of polynomial order $3$ and
window $\min(301,\,2\lfloor n/4\rfloor+1)$ cycles for a cell of $n$ cycles,
and the derivative is divided by the standard deviation of that channel's
residual about the same filter, which whitens the channel by its own measured
noise. The whitened derivative is interpolated onto $40$ points of
$\tau\in[0.15,0.90]$. One sign per channel, the median sign over every cell
and grid point, orients all cells alike; each cell's vector is then
normalised at every grid point, and the fleet mean, renormalised, is the
empirical informative direction. Its arc is the sum of the angles between
consecutive grid points, and its end-to-end angle is the angle between the
first and the last. Positivity is the fraction of grid points at which every
component of the fleet mean is positive; it is one at the window used.

\subsection{Why three measured fleets cannot test actuation}

Testing actuation needs the operating point to change within a unit while
everything else is held. In XJTU-SY each of three operating conditions is
applied to its own five bearings, and in PRONOSTIA each of three conditions
to its own bearings, so a difference between conditions cannot be told apart
from a difference between groups of units. The NASA battery fleet has four
cells. These fleets are therefore listed with the reason and no statistic.

\subsection{The actuation test on C-MAPSS}

The training files of FD004 and FD002 are used. Life is normalised per engine
by its final cycle. Operating regimes are labelled by standardising the three
operating settings over the file and rounding to one decimal; each file has
$6$ regimes. A sensor is kept if its standard deviation within a
(engine, regime) group has a positive median, and each kept sensor is
standardised within each regime, so that a regime's level does not enter the
direction.

Within one engine, each regime's cycles are resampled to a common count $m$,
the smallest over that engine's regimes, and an engine enters only if
$m\ge2(w+3)$. Each resampled series is split into its even and its odd
cycles, two half-samples of the same regime. On each half the direction is
estimated as for the Severson cells, with polynomial order $2$ and window $w$,
and interpolated onto $10$ points of $\tau\in[0.30,0.85]$. The angle between
two directions is the mean over grid points of the pointwise angle. $W$ is the
mean, over the engine's regimes, of the angle between the two halves of the
same regime, and $B$ the mean, over pairs of different regimes, of the angle
between a half of one and a half of the other. $B-W$ is positive when regimes
differ by more than half-samples of one regime differ from each other, which
is what an operating point that moves the direction predicts. The table
reports the mean of $B-W$ over engines, and the number of engines for which it
is positive, at $w=7$; at $w=9$ and $w=11$ every value keeps its sign and
exceeds the largest control bias.

The negative control uses the same machinery on fake regimes. Within one
regime of one engine with at least $4(w+3)$ cycles, the four interleaved
subsequences of every fourth cycle stand in for four regimes; $W$ compares the
first with the second and $B$ the first with the fourth. No contrast exists by
construction. Over all three windows and both files the mean of $B-W$ stays
within $0.5^{\circ}$ of zero, and a paired $t$ test of $B$ against $W$ does
not reject equality at any window, the smallest $p$ being $0.35$.

\subsection{The separability test on N-CMAPSS}

The development sets of DS02, units $2$, $5$, $10$, $16$, $18$ and $20$, and of
DS01, units $1$ to $6$, are used. The operating point is held fixed by
clustering flight points, by $k$-means with $6$ clusters, on standardised
altitude, Mach number and throttle-resolver angle; life is divided into $2$
stages of equal width in cycles. A (stage, cluster) cell enters only with at
least $400$ points.

Within a cell each sensor is regressed on standardised remaining life, and
the slope is divided by the standard deviation of that sensor's residual, so
that the vector of slopes is the informative direction at a fixed operating
point rather than a residual after standardisation. The vector is normalised,
and the logarithms of the absolute values of its components, centred on
their mean, are what the test compares. The logarithm is the point: under the
separable form of Lemma~\ref{M-lem:sep} each component is a product of an
operating-point factor and a state factor, so in logarithms the two act
additively, and separability is the absence of interaction.

Each cell is estimated twice, from its even and from its odd points. The mean
of the two is the signal; their half-difference measures noise, and applying
the same decomposition to it gives each effect its own noise floor. The
decomposition is the ordinary two-way one over the cells, into a grand mean,
a stage effect, an operating-point effect and an interaction, with sums of
squares taken over cells and sensors and pooled over units. With $6$ clusters
and $2$ stages the operating-point effect and the interaction each have $5$
degrees of freedom per sensor and the stage effect $1$, so the reported ratio
of interaction to operating-point effect compares quantities with equal
degrees of freedom. That ratio is $1.65$ for DS02 and $1.84$ for DS01. Against
its own noise floor the interaction is $566$ times larger on DS02
and $110$ times on DS01, so it is not noise. The test was repeated
at $3$, $4$, $6$, $8$ and $10$ clusters with $2$ stages, at $6$ clusters with
$3$ stages and at $4$ clusters with $4$ stages, with at least $300$ points per
cell; the interaction exceeds the operating-point effect in every
configuration, by a factor of at least $1.42$.

\begin{table}[tb]
\caption{The parameters of each test as run.}
\label{tab:params}
\centering
\footnotesize
\setlength{\tabcolsep}{3pt}
\begin{tabular}{@{}>{\raggedright\arraybackslash}p{0.17\textwidth}
                  >{\raggedright\arraybackslash}p{0.78\textwidth}@{}}
\hline
test & parameters\\
\hline
Severson rotation & $129$ cells with at least $400$ cycles; Savitzky--Golay
order $3$, window $301$ cycles or $2\lfloor n/4\rfloor+1$ if smaller; $40$
points on $\tau\in[0.15,0.90]$\\[2pt]
C-MAPSS actuation & FD004 and FD002 training sets; $6$ regimes each; order
$2$, window $w=7$, stability at $w=9,11$; $10$ points on
$\tau\in[0.30,0.85]$; engine enters if $m\ge2(w+3)$; control needs
$4(w+3)$ cycles in one regime\\[2pt]
N-CMAPSS separability & DS02 units $2,5,10,16,18,20$ and DS01 units $1$--$6$;
$6$ clusters, $2$ stages, at least $400$ points per cell; sweep over seven
(cluster, stage) settings with at least $300$ points\\
\hline
\end{tabular}
\end{table}

\section{Proofs of Propositions~\ref{M-prop:record} and~\ref{M-prop:bary}}
\label{sec:proofs}

\begin{proof}[Proof of Proposition~\ref{M-prop:record}]
In whitened coordinates the linearised record is $\tilde y_k=\dvec_k R+\epsilon_k$ with $\epsilon_k\sim N(0,I_d)$ independent, so the Fisher information it carries about $R$ is $\sum_k\dvec_k^{\top}\dvec_k=\sum_k\|\dvec_k\|^{2}$ and an oracle reading the whole record obtains $\tfrac12\log(1+G)$ nats. The fixed indicator returns the scalars $z_k=\mathbf v^{\top}\tilde y_k=(\mathbf v^{\top}\dvec_k)R+\mathbf v^{\top}\epsilon_k$, whose noise terms $\mathbf v^{\top}\epsilon_k$ are standard normal and independent because $\|\mathbf v\|=1$, so the compressed record carries Fisher information $\sum_k(\mathbf v^{\top}\dvec_k)^{2}$. Since $\mathbf v^{\top}\dvec_k=\|\dvec_k\|\cos\psi_k$,
\[
\Var[R]\sum_k(\mathbf v^{\top}\dvec_k)^{2}
=\Var[R]\sum_k\|\dvec_k\|^{2}\cos^{2}\psi_k=G(1-S),
\]
and subtracting the two mutual informations gives $\tfrac12\log(1+G)-\tfrac12\log\big(1+G(1-S)\big)=-\tfrac12\log\frac{1+G-GS}{1+G}=-\tfrac12\log(1-WS)$.
\end{proof}

\begin{proof}[Proof of Proposition~\ref{M-prop:bary}]
(i) As $\gamma\to\infty$ we have $w\to1$, and \eqref{M-eq:exact} gives
\[
\lloss\to-\tfrac12\log\cos^{2}\psi=-\log\sum_i\sqrt{p_iq_i}=\DB
\]
by Proposition~\ref{M-prop:bhat}. The convergence is uniform, which is what
carries the minimisers with it. Positivity keeps $p_\tau$ interior, and
$\sqrt{q_i}\ge q_i$ with $\sum_iq_i=1$ gives
$\cos\psi=\sum_i\sqrt{p_iq_i}\ge\min_i\sqrt{p_i}$ for every $q$ in the
simplex, so $\cos\psi$ is bounded away from zero uniformly in $q$ and in
$\tau$ and $\DB$ stays finite. Both objectives are continuous on a compact set
and uniformly close on it, so their minimisers converge. (ii) By
\eqref{M-eq:expansion}, $\bar\lloss=\tfrac12\bar
w\,\E_\mu[\psi^{2}]+O(\psi^{4})$ with $d\mu=w\,d\tau/\bar w$; minimising a
$\mu$-weighted mean square of a geodesic distance is the definition of the
weighted Fr\'echet mean, and $\psi=\tfrac12\dFR$ converts $\tfrac12\bar
w\,\E_\mu[\psi^{2}]$ into $\tfrac18\bar w\,\E_\mu[\dFR^{2}]$.
\end{proof}

\section{Two candidate certificates for Proposition~\ref{M-prop:converse}}
\label{sec:certificates}

The $C^{1}$ hypothesis cannot be weakened to Lipschitz continuity. A Lipschitz
$V$ is differentiable almost everywhere in the state by Rademacher's theorem,
but a single trajectory is a null set and could in principle run inside the
exceptional set, so the chain rule step would be unavailable; recovering the
argument at that generality needs the Clarke generalised gradient in place of
$\nabla V$. Smoothness is also, as it turns out, exactly what the natural
candidates lack.

We tried two. The value function of the Dinkelbach problem of Section~\ref{M-sec:limit} satisfies
\eqref{M-eq:cert} with equality in the viscosity sense~\cite{bardi1997}, but
value functions of this class are only Lipschitz. Computed on a $41^{3}$ grid
and tested on a $61^{3}$ one, \eqref{M-eq:cert} fails at about $45\%$ of test
points by up to $2.9\times10^{-3}$. Halving $\lambda$ moves that violation by
$0.6\%$, so the fault lies in the candidate and not in the bound. At the
opposite extreme an affine $V(x)=c_0-a_0^{\top}x$ has an exactly constant
gradient and makes every requirement linear, so the best certifiable $\lambda$
is a linear program; over $167{,}936$ constraints it returns $V\equiv0$ and
only the pointwise bound $\min_{x,u}\lloss$, which is vacuous here because the
reachable floor falls to within sampling resolution of zero at mid-life.
Deleting the two boundary requirements lets that program run to any cap
imposed on it, which identifies them as what binds: trajectories of
\eqref{M-eq:plant} increase in every coordinate, so an affine $V$ is monotone
along all of them at once and $V(x_0)\ge0\ge V|_{\text{failure}}$ forces it to
decrease, while the near-aligned states, where the running cost vanishes,
demand a local increase.

The two failures are complementary, and together they say what a certificate
needs. The value function varies its gradient exactly as the two demands
require but is not differentiable where that matters; affine functions are
perfectly smooth but cannot vary their gradient at all. A certificate must
therefore be both smooth and non-constant in its gradient, which points at a
smooth under-approximation of the value function, obtained by sum-of-squares
programming when the plant data are polynomial, or by interval arithmetic with
an explicit Lipschitz constant.

\section{Numerical verification}
\label{sec:verification}

\subsection{The identities}

The analytical results of the main text were checked numerically before
they were used. The rotation identity \eqref{M-eq:repl} holds to
relative error below $1.1\times10^{-7}$ over the five input--channel pairs
$(d,m)=(2,1)$, $(3,2)$, $(3,1)$, $(4,2)$ and $(5,3)$; the replicator form to
$6.9\times10^{-8}$; Theorem~\ref{M-thm:beta} to $5.1\times10^{-6}$, limited by
finite differencing; and \eqref{M-eq:bhat} to $1.1\times10^{-16}$.

\subsection{The exceptional family of Remark~\ref{M-rem:family}}

Integrating the construction of Remark~\ref{M-rem:family}
for $d=2$ to $8$ with one input holds the misalignment below
$1.5\times10^{-6}$ degrees and the mean loss below $10^{-15}$. These families carry no prior variance of their own, so their losses are evaluated at a fixed $w=0.9$ rather than at a $\gamma$ the plant supplies; the bound has fifteen orders of magnitude of headroom and does not turn on the choice. That is against
$1.2\times10^{-5}$ to $5.4\times10^{-4}$ for the best constant input on the
same plant.

The family is a knife edge. Perturbing the shape coefficients $\nu$ by $1\%$ raises the mean loss to
$3.2\times10^{-4}$ under the myopic policy, against $1.8\times10^{-4}$ for the
best constant input on the same perturbed plant. At one percent of shape
error, steering is already worse than not steering. Closing the loop halves
the damage but does not repair it, so the fragility belongs to the
construction and not to the open-loop schedule. Mistuning the activation
energies instead is tolerated to about two percent.

\bibliographystyle{elsarticle-num}
\bibliography{Ref}

\end{document}
